\documentclass[12pt]{article}
 
% ── Packages ──────────────────────────────────────────────────────────────────
\usepackage{amsmath, amssymb, amsthm}
\usepackage{geometry}
\usepackage{booktabs}
\usepackage{array}
\usepackage{xcolor}
\usepackage{hyperref}
 
\geometry{margin=1in}
 
% ── Theorem environments ───────────────────────────────────────────────────────
\newtheorem*{problem}{Problem}
\newtheorem*{claim}{Claim}
\theoremstyle{definition}
\newtheorem*{example}{Example}
\theoremstyle{remark}
\newtheorem*{remark}{Remark}
 
% ── Custom commands ────────────────────────────────────────────────────────────
\newcommand{\R}{\mathbb{R}}
\newcommand{\nullT}{\operatorname{null}}
\newcommand{\range}{\operatorname{range}}
\newcommand{\dime}{\operatorname{dim}}
 
% ── Document ───────────────────────────────────────────────────────────────────
\begin{document}
 
\begin{center}
  {\Large \textbf{Linear Algebra Done Right} (4th ed.) --- Sheldon Axler}\\[6pt]
  {\large Exercises 3B, Problem 1}\\[4pt]
  {\normalsize USF Master's Preparation}
\end{center}
 
\hrule
\vspace{1em}
 
% ── Problem statement ──────────────────────────────────────────────────────────
\begin{problem}
Give an example of a linear map $T$ with $\dime\,\nullT\,T = 3$ and
$\dime\,\range\,T = 2$.
\end{problem}
 
\vspace{0.5em}
 
% ── Background theorem ─────────────────────────────────────────────────────────
\noindent\textbf{Key Tool: Fundamental Theorem of Linear Maps.}\;
For any linear map $T : V \to W$ with $V$ finite-dimensional,
\[
  \dime\,V \;=\; \dime\,\nullT\,T \;+\; \dime\,\range\,T.
\]
Both examples below satisfy $3 + 2 = 5$, confirming that the domain must have
dimension at least $5$.
 
\vspace{1em}
\hrule
\vspace{1em}
 
% ── Example 1 ─────────────────────────────────────────────────────────────────
\begin{example}[Same domain and codomain]
Define $T : \R^5 \to \R^5$ by
\[
  T(x_1, x_2, x_3, x_4, x_5) \;=\; (x_1,\, x_2,\, 0,\, 0,\, 0).
\]
\end{example}
 
\begin{claim}
$\dime\,\nullT\,T = 3$ and $\dime\,\range\,T = 2$.
\end{claim}
 
\begin{proof}
\textbf{Null space.}
A vector $(x_1,\ldots,x_5)\in\R^5$ lies in $\nullT\,T$ if and only if
$T(x_1,\ldots,x_5) = \mathbf{0}$, i.e.\ $x_1 = x_2 = 0$ while $x_3,x_4,x_5$
are free.  Hence
\[
  \nullT\,T
  = \{(0,\,0,\,x_3,\,x_4,\,x_5) : x_3,x_4,x_5\in\R\}
  = \operatorname{span}\!\bigl(\,e_3,\,e_4,\,e_5\bigr),
\]
so $\dime\,\nullT\,T = 3$.
 
\textbf{Range.}
Every output has the form $(x_1, x_2, 0, 0, 0)$, so
\[
  \range\,T
  = \{(x_1,\,x_2,\,0,\,0,\,0) : x_1,x_2\in\R\}
  = \operatorname{span}\!\bigl(\,e_1,\,e_2\bigr),
\]
giving $\dime\,\range\,T = 2$.
Note that $\range\,T$ is a \emph{proper subspace} of the codomain $\R^5$.
 
\textbf{Sanity check.}
$\dime\,\R^5 = 5 = 3 + 2 = \dime\,\nullT\,T + \dime\,\range\,T$.\;\checkmark
\end{proof}
 
\vspace{1em}
\hrule
\vspace{1em}
 
% ── Example 2 ─────────────────────────────────────────────────────────────────
\begin{example}[Distinct domain and codomain]
Define $T : \R^5 \to \R^2$ by
\[
  T(x_1, x_2, x_3, x_4, x_5) \;=\; (x_1,\, x_2).
\]
\end{example}
 
\begin{claim}
$\dime\,\nullT\,T = 3$ and $\dime\,\range\,T = 2$.
\end{claim}
 
\begin{proof}
\textbf{Null space.}
$T(x_1,\ldots,x_5) = \mathbf{0}\in\R^2$ if and only if $x_1 = x_2 = 0$, with
$x_3,x_4,x_5$ free.  Thus $\nullT\,T = \operatorname{span}(e_3,e_4,e_5)$ and
$\dime\,\nullT\,T = 3$.
 
\textbf{Range.}
For any $(y_1,y_2)\in\R^2$, we have $T(y_1,y_2,0,0,0)=(y_1,y_2)$, so $T$ is
\emph{surjective} and $\range\,T = \R^2$, giving $\dime\,\range\,T = 2$.
 
\textbf{Sanity check.}
$\dime\,\R^5 = 5 = 3 + 2$.\;\checkmark
\end{proof}
 
\vspace{1em}
\hrule
\vspace{1em}
 
% ── Comparison table ───────────────────────────────────────────────────────────
\subsection*{Comparison of the Two Examples}
 
\begin{center}
\renewcommand{\arraystretch}{1.4}
\begin{tabular}{@{} l >{\centering\arraybackslash}m{3.8cm}
                      >{\centering\arraybackslash}m{3.8cm} @{}}
\toprule
 & \textbf{Example 1} & \textbf{Example 2} \\
\midrule
Map & $T:\R^5\to\R^5$ & $T:\R^5\to\R^2$ \\
Formula & $(x_1,x_2,x_3,x_4,x_5)\mapsto(x_1,x_2,0,0,0)$ & $(x_1,x_2,x_3,x_4,x_5)\mapsto(x_1,x_2)$ \\
$\dime\,\nullT\,T$ & $3$ & $3$ \\
$\dime\,\range\,T$ & $2$ & $2$ \\
Is $T$ surjective? & No \;($\range\,T\subsetneq\R^5$) & \textbf{Yes} \;($\range\,T = \R^2$) \\
Codomain $=$ Range? & No & \textbf{Yes} \\
\bottomrule
\end{tabular}
\end{center}
 
\begin{remark}
The key distinction is that Example~1 maps into a \emph{larger} codomain
$\R^5$, so the range is a proper $2$-dimensional subspace of $\R^5$ and $T$ is
not surjective.  Example~2 maps into $\R^2$, which exactly equals the range,
making $T$ surjective. Both are valid answers to the problem, but
Example~2 is arguably more \emph{minimal}: the codomain is chosen to match the
desired range dimension exactly, leaving no ``wasted'' space.
\end{remark}
 
\end{document}